\documentclass[11pt]{article}
\usepackage[margin=1in]{geometry}
\usepackage{graphicx,booktabs,float,xcolor,hyperref,amssymb}
\usepackage{helvet}\renewcommand{\familydefault}{\sfdefault}
\hypersetup{colorlinks=true,linkcolor=blue,urlcolor=blue}
\newcommand{\Mw}{$M_{\mathrm w}$}
\setlength{\parindent}{0pt}\setlength{\parskip}{6pt}
\title{\textbf{Can teleseismic data resolve a finite-fault model of an \Mw~6.1?}\\
\large The 2023 Mariana Islands earthquake (us6000kzv3) as a resolution test of the
MTUQ$+$WISP pipeline}
\author{Automated pipeline (17.Venezuela\_2026)}
\date{2026-07-29}
\begin{document}\maketitle

\section*{Purpose}
USGS publishes finite-fault models routinely for \Mw$\gtrsim$6.5; for the 2023 \Mw~6.1 Mariana
Islands earthquake no finite-fault product exists. Hartzell, Mendoza \& Zeng (2013, GRL)
nevertheless resolved the \Mw~5.8 Mineral, Virginia rupture from teleseismic P waves alone ---
their enabling ingredients being P waves to $\sim$1 Hz (attenuation $t^*\!\approx\!1$ s still
passes source detail at teleseismic range), resolution carried by \emph{timing} rather than
amplitude, km-scale subfaults, and a compact fault. This report applies that recipe with the
MTUQ$+$WISP pipeline and states explicitly what is, and is not, resolved.

\section{Event and reference solution}
\begin{table}[H]\centering\small
\begin{tabular}{ll}
\toprule
Origin & 2023-08-14T13:51:53.8 UTC, 13.347$^\circ$N, 147.530$^\circ$E, depth 8 km \\
Magnitude & \Mw~6.1 (W-phase); $M_0=1.74\times10^{18}$ N\,m; centroid 13.5 km; DC 91\,\% \\
Tectonics & oblique normal $+$ strike-slip east of Guam, near the trench (outer-rise style) \\
W-phase NP1 & 253.3/56.1/$-$28.7 \\
W-phase NP2 & 000.3/66.5/$-$142.6 \\
USGS products & MT, ShakeMap, PAGER, DYFI --- \textbf{no finite fault} \\
\bottomrule
\end{tabular}
\end{table}

\section{Data}
205 three-component stations, 32--89$^\circ$, maximum azimuthal gap 28$^\circ$ --- the densest
coverage of any event in this project, which matters doubly here: the mechanism has a large
strike-slip component, and the Kumamoto blind test showed clustered azimuths alias strike-slip
radiation. MTUQ used an azimuth-balanced 60-station response-corrected subset; WISP applied its
own quality control. WISP's body-wave filter is 0.01--1.0 Hz --- the Hartzell band --- with
auto-sized $2.4\times2.0$ km subfaults (255 total).

\begin{figure}[H]\centering
\includegraphics[width=0.62\textwidth]{../fig_mariana_stations.png}
\caption{Teleseismic stations (azimuthal-equidistant; rings 30/60/90$^\circ$).}
\end{figure}

\section{Moment tensor (MTUQ): the independent point source fails at \Mw~6.1}
At $M_0\sim2\times10^{18}$ N\,m the standard teleseismic bands (body 10--33 s, surface
50--100 s) are low-SNR at 30--90$^\circ$, and it shows:

\begin{table}[H]\centering\small
\begin{tabular}{llll}
\toprule
variant & mechanism & \Mw & Kagan vs W-phase \\
\midrule
DC joint, 60 stations & 68/82/$-$2 & 6.25 & 52.5$^\circ$ \\
DC body P$+$SH, 60 stations & 122/54/52 & 5.80 (lower grid edge) & 85.7$^\circ$ \\
depth scan 6--33 km & mechanism drifts & --- & misfit varies 0.2\,\% \\
\bottomrule
\end{tabular}
\caption{Neither the mechanism nor the centroid depth is independently resolvable; the
depth-misfit curve is flat to 0.2\,\% over 6--33 km. This is not an azimuth-coverage problem
(gap 28$^\circ$) --- it is a signal-to-noise floor. Pipeline rule for \Mw$\lesssim$6.5: take
mechanism and depth from the W-phase and spend the teleseismic data on the finite fault ---
exactly the Hartzell strategy.}
\end{table}

\begin{figure}[H]\centering
\includegraphics[width=0.8\textwidth]{../fig_mariana_mtuq_beachballs.png}
\caption{W-phase reference vs the two MTUQ attempts.}
\end{figure}

\subsection*{Why the MTUQ result is poor: the waveforms show it directly}
The fits below make the failure mode concrete. Observed peak amplitudes are \textbf{1--5 nm}
(body) and a few mm (long-period surface waves) --- at or barely above the noise floor of the
teleseismic bands. Most of the $\sim$180 traces are near-flat; a handful carry large
oscillations that are noise bursts, not signal, and the normalized L2 misfit is then steered by
those few traces (the low-effective-trace-count failure familiar from the regional Korea
pipeline). Crucially, evaluating the \emph{W-phase mechanism itself} on the same processed data
(second figure) fits the usable traces essentially as well as the grid-search best --- the
misfit surface is flat between mechanisms 50$^\circ$ apart, i.e.\ the data do not
discriminate. The poor MTUQ result is a property of the magnitude, not of the protocol.

\begin{figure}[H]\centering
\includegraphics[width=0.5\textwidth]{../../results/mtuq_dc_2023mariana/waveforms_joint.png}
\caption{MTUQ joint DC fits at the grid-search best (68/82/$-$2). Black observed, red
synthetic; most traces are near-flat at nm level.}
\end{figure}
\begin{figure}[H]\centering
\includegraphics[width=0.5\textwidth]{../../results/mtuq_dc_2023mariana/2023mariana_waveforms_wphase.png}
\caption{The same data fit with the mechanism FIXED to the USGS W-phase solution
(253/56/$-$29, \Mw~6.1): visually indistinguishable fit quality from the grid best ---
direct evidence that the teleseismic bands cannot discriminate the mechanism at this
magnitude.}
\end{figure}

\section{Finite fault (WISP): both W-phase planes, Mww-seeded, \texttt{shift\_match} applied}
\begin{table}[H]\centering\small
\begin{tabular}{lllll}
\toprule
plane & strike/dip/rake & misfit (post-shift) & peak slip & T$_{5-95}$ \\
\midrule
\textbf{NP1 (preferred)} & \textbf{253/56/$-$29} & \textbf{0.2012} & 24.9 cm & 8.2 s \\
NP2 & 0/66/$-$143 & 0.2058 & 24.9 cm & 7.9 s \\
\bottomrule
\end{tabular}
\caption{Recovered $M_0=1.76\times10^{18}$ N\,m vs W-phase $1.74\times10^{18}$ --- 1\,\%
agreement. The 2.2\,\% NP1 preference sits between the Calama near-tie (1.6\,\%, $\sim$10 km
rupture) and the SPdA decisive margin (10\,\%, $\sim$50 km rupture): suggestive, not
conclusive, exactly as the rupture-size criterion predicts.}
\end{table}

\begin{figure}[H]\centering
\includegraphics[width=0.8\textwidth]{../20230814135153/ffm.0/NP1_shift/plots/SlipIsochrone_trueScale.png}
\caption{Preferred plane (253/56/$-$29), true 1:1 scale: a single compact asperity
$\sim$10 km across at 2--13 km depth, centred on the 8-km hypocentre; peak slip 25 cm;
concentric 1-s isochrones spanning 0--5 s (V$_r\approx2.0$ km s$^{-1}$, subshear); oblique
right-lateral$+$normal slip vectors.}
\end{figure}
\begin{figure}[H]\centering
\includegraphics[width=0.55\textwidth]{../20230814135153/ffm.0/NP1_shift/plots/MomentRate.png}
\caption{Moment-rate function (preferred plane).}
\end{figure}
\begin{figure}[H]\centering
\includegraphics[width=0.92\textwidth]{../20230814135153/ffm.0/NP1_shift/plots/Map.png}
\caption{Map view.}
\end{figure}
\begin{figure}[H]\centering
\includegraphics[width=0.92\textwidth]{../20230814135153/ffm.0/NP1_shift/plots/P_body_waves.png}
\caption{Teleseismic P fits (preferred plane, shifted) --- the dataset that carries the
Hartzell-style timing resolution.}
\end{figure}
\begin{figure}[H]\centering
\includegraphics[width=0.92\textwidth]{../20230814135153/ffm.0/NP1_shift/plots/SH_body_waves.png}
\caption{Teleseismic SH fits (preferred plane, shifted).}
\end{figure}
\begin{figure}[H]\centering
\includegraphics[width=0.92\textwidth]{../20230814135153/ffm.0/NP1_shift/plots/Rayleigh_surf_waves.png}
\caption{Rayleigh fits. WISP's quality control retained only five usable long-period vertical
channels --- mm-scale mantle waves sit at the noise floor for \Mw~6.1 --- but the retained
traces fit well (SYDS is a zero-weight rejected trace shown for completeness). Surface waves
constrain moment, not slip detail, at this magnitude.}
\end{figure}
\begin{figure}[H]\centering
\includegraphics[width=0.92\textwidth]{../20230814135153/ffm.0/NP1_shift/plots/Love_surf_waves.png}
\caption{Love fits (same caveat).}
\end{figure}

\subsection*{Why the body-wave window after $\sim$15 s is rarely fit}
A systematic feature of the P fits: the direct pulse (0--10 s, including pP/sP at 4.5--6.5 s
for the 13.5-km source) is matched at nearly every station, while observed energy after
$\sim$15 s is left unfit. Two causes, both identifiable:
\begin{enumerate}
\item \textbf{The velocity model has no ocean.} The source-side model (table below) tops out
in a 90-m soft-sediment layer --- but this trench-adjacent source lies beneath $\sim$4--5 km
of Pacific water. Water-column multiples (pwP and its reverberations, two-way time
$2\times4.5/1.5\approx6$ s per bounce) form trains persisting tens of seconds that a
land-style 1-D model \emph{cannot produce at all}. This is a known limitation of the
operational WISP crust model for oceanic events, not a wrong crustal layer.
\item \textbf{The oceanic microseism band.} Several stations show the same 5--8-s
oscillations \emph{before} the P arrival at the same nm amplitude --- the double-frequency
microseism peak lives at exactly these periods and amplitudes, so part of the late window is
simply noise that no source model should fit.
\end{enumerate}
\begin{table}[H]\centering\small
\begin{tabular}{lllll}
\toprule
$V_p$ (km/s) & $V_s$ (km/s) & $\rho$ (g/cm$^3$) & thickness (km) & \\
\midrule
2.56 & 0.60 & 1.82 & 0.09 & (soft sediment --- no water layer above) \\
5.37 & 2.97 & 2.42 & 1.12 & \\
6.45 & 3.67 & 2.71 & 2.45 & \\
7.36 & 4.20 & 2.90 & 7.57 & \\
8.08 & 4.47 & 3.38 & 196 & \\
\bottomrule
\end{tabular}
\caption{WISP source-side velocity model. The absence of a water column is the structural
reason the post-15-s window cannot be fit.}
\end{table}
The practical consequence is benign for the headline conclusions --- the inversion is anchored
by the direct-pulse \emph{timing} (the Hartzell resolution carrier), and the plane-stable STF
confirms it --- but it inflates the absolute misfit (0.20 vs 0.15 for continental Kumamoto)
and is the first thing to fix (add a water layer / use an oceanic crust profile) if slip
detail beyond first order is ever needed for a trench event.

\section{Resolution assessment}
\begin{figure}[H]\centering
\includegraphics[width=\textwidth]{../fig_mariana_resolution.png}
\caption{Left: the source-time function is stable across the two (very different) assumed
planes, $r=0.98$, including its two-pulse structure at 2.8 and 5.4 s --- the plane-independent,
therefore data-driven, part of the model. The Ekstr\"om-scaling duration band ($2\cdot
h_{dur}=5.8$ s) captures the STF body; the tail beyond it is regularization. Centre: the MTUQ
point-source depth scan is flat to 0.2\,\% --- depth must come from the W-phase. Right: the
NP1-vs-NP2 misfit margin in the context of the pipeline's plane-discrimination criterion.}
\end{figure}

\textbf{Resolved} (stable across planes, passing physical checks):
\begin{itemize}
\item a single compact asperity $\sim$10 km across at 2--13 km depth, centred on the
hypocentre --- consistent with \Mw~6.1 source scaling;
\item the STF shape including its two-pulse structure (2.8 $+$ 5.4 s), cross-plane $r=0.98$;
\item $M_0$ to 1\,\% of the W-phase; rupture front spanning 0--5 s, V$_r\approx2.0$ km s$^{-1}$;
\item stress drop 1.0 MPa (effective-area convention, $(\Sigma u)^2/\Sigma u^2 \to 263$ km$^2$)
to 3.7 MPa (half-peak area, 110 km$^2$) --- in the physically expected 1--10 MPa range.
\end{itemize}

\textbf{Not resolved} (and stated so):
\begin{itemize}
\item the fault plane: the 2.2\,\% NP1 margin is suggestive, not decisive, at 10-km rupture size;
\item an independent point-source mechanism or centroid depth (Section 3);
\item fine slip detail and the duration tail: T$_{5-95}=8.2$ s exceeds the 5.8-s scaling
expectation because smoothing spreads low-level slip; the 25-cm peak is a lower bound;
\item absolute updip/downdip placement within the compact patch.
\end{itemize}

\section{Verdict}
\textbf{Yes: a reliable first-order finite-fault model of an \Mw~6.1 is recoverable from
teleseismic data} --- when the Hartzell recipe is followed (W-phase-fixed mechanism, P waves to
1 Hz, km-scale subfaults, claims restricted to plane-stable features). What cannot be claimed
at this magnitude: the fault plane, an independent mechanism, or a centroid depth from
teleseismic point-source inversion. This extends the pipeline's validated range from \Mw~6.8
down to \Mw~6.1, one magnitude unit above the Hartzell et al.\ (2013) \Mw~5.8 benchmark.

\end{document}
